\documentclass[a4paper, oneside]{article}
\usepackage{graphicx}
\usepackage{amsthm}
\usepackage{amsmath}
\usepackage{amssymb}
\usepackage[a4paper,
            bindingoffset=0.2in,
            left=2cm,
            right=2cm,
            top=2cm,
            bottom=2cm,
            footskip=.25in]{geometry}
\usepackage[italian]{babel}
\usepackage{pgfplots}
\usepackage{tabularx}
\usepackage{tikz}
\usepackage{wrapfig}
\usepackage{color}
\usepackage[d]{esvect}
\definecolor{page}{rgb}{0.129,0.157,0.212}
\pagecolor{page}
\color{white}
\graphicspath{ {./images/} }
\usetikzlibrary{shapes.geometric}
\usetikzlibrary{datavisualization}
\usetikzlibrary{datavisualization.formats.functions}
\usetikzlibrary{patterns}
\pgfplotsset{width=10cm,compat=1.18}

\title{Appunti di Fisica II}
\author{Tommaso Miliani}
\date{09-03-26}

\begin{document}
\newtheoremstyle{theoremEnv}
                {}          % Space above
                {}          % Space below
                {\slshape}  % Body font
                {}          % Indent amount
                {\bfseries} % Head font
                {.}         % Punctuation after head
                {\newline}  % Space after theorem head
                {}          % Theorem head spec
\theoremstyle{theoremEnv}

\newtheorem{definition}{Definizione}[section]
\newtheorem{theorem}{Teorema}[section]
\newtheorem{lemma}{Proposizione}[section]
\newtheorem{observation}{Osservazione}[section]
\newtheorem{corollary}{Corollario}[theorem]
\newtheorem{example}{Esempio}[section]
\newtheorem{remark}{Enunciato}[section]

\maketitle

\section{Conduttore in statica}
Per noi i conduttori sono dei materiali in cui esistono un numero illimitato
di cariche che comporta che la forza tra le cariche sia nulla: il campo interno ad un conduttore,
qualsiasi forma abbia, è nullo. Questo vuol dire anche che il conduttore è un 
volume equipotenziale: se il potenziale variasse, allora si avrebbe campo. All'interno 
del quale la densità di carica è nulla: le cariche positive sono fisse mentre quelle 
negative sono libere di muoversi, facendo sì che per ogni volumetto 
la carica totale sia nulla, anche se sulla superficie esiste una certa densità di carica.
\section{Capacità di un conduttore}
\begin{wrapfigure}{r}{0.4\textwidth}
    \centering
    \caption{Il conduttore sferico}
    \begin{tikzpicture}
        \draw(0, 0) circle (1);
        \draw(0, 0) -- (0.5, 0.83) node[midway, left] {$R$};
        \node at (0.7, -1) {$Q$};
    \end{tikzpicture}    
\end{wrapfigure}
Considerata una sfera conduttrice sulla cui superficie esterna si ha 
un eccesso di carica positiva (per cui si può estrarre degli elettroni 
dal conduttore). Se il problema è a simmetria sferica $\sigma$ non dipende
né da $\phi$ né da $\theta$, dunque per una sfera di raggio $R$: 
\begin{gather*}
    \sigma = \frac{Q}{4\pi R^{2}}
\end{gather*}
Allora il campo elettrico della sfera, data la simmetria del problema:
\begin{gather*}
    \vv{E}(r, \theta, \phi) = E_r(r) \hat{r} \ \Longrightarrow \ E_r = \left\{\begin{array}{l}
        0 \qquad\quad r < R \\
        \dfrac{q}{4\pi\epsilon_0 r^{2}} \ r > R
    \end{array}\right.
\end{gather*}
Sulla superficie della sfera il campo elettrico è discontinuo, sulla superficie di 
un conduttore di qualunque forama si torna al teorema di Coulomb per cui 
\begin{gather*}
    E_r = \frac{\sigma}{\epsilon_0} \qquad r = R
\end{gather*}
Per poter arrivare al concetto di Capacità di deve osservare il potenziale del conduttore 
all'infinito. In generale si può esprimere il potenziale come 
\begin{gather*}
    V(\infty ) - V(R) = -\int_{R}^{\infty } \vv{E} \cdot \vv{dr} = -\frac{Q}{4\pi\epsilon_0} \int_{R}^{\infty } \frac{1}{r^{2}} dr = -\frac{Q}{4\pi\epsilon_0} \frac{1}{R}
\end{gather*}
Utilizzando la convenzione secondo la quale $V(\infty ) = 0$, si ottiene che il potenziale 
della sfera è esattamente dato da
\begin{align}
    V(R) = \frac{Q}{4\pi\epsilon_0 R} \equiv \frac{Q}{\mathcal{C}} \ \Longrightarrow \ \mathcal{C} = 4\pi\epsilon_0 R
\end{align}
La cui unità di misura è "Farad", ossia Coulomb su metro.
\begin{example}
    Una sfera metallica di raggio un metro ha una 
capacità di $10^{-10}$ F, per cui se la si portasse a $10^{2}$ V, essa avrebbe solo $10^{-8}$ Coulomb. Nonostante
non sia molta carica, si avrebbe una forza pari a
\begin{gather*}
    F = \frac{1}{4\pi\epsilon_0} \frac{q^{2}}{R^{2}} = 10^{-2} N
\end{gather*}
Ossia una forza molto piccola. I metalli non hanno una grande capacità: il problema è accumulare tanta 
carica su di un dispositivo in modo tale che non abbia una differenza di potenziale troppo alto.
\end{example} 

\section{Gli accumulatori di carica: i condensatori}
\subsection{Condensatori sferici}
\begin{wrapfigure}{r}{0.4\textwidth}
    \centering
    \caption{Il condensatore sferico con la sfera esterna messa a terra.}
    \begin{tikzpicture}
        \draw(0, 0) circle (1);
        \draw(0, 0) circle (1.45);
        \draw(0, 0) circle (1.5);
        \draw(0, 0) -- (1, 0) node[midway, above] {$R_1$};
        \draw(0, 0) -- (0, 1.5) node[midway, right] {$R_2$};
        \node at (0, -1.2) {$Q$};
        \node at (1.3, 1.3) {$-Q$}; 
        \draw(1.5, 0) -- (1.7, 0) -- (1.7, -1);
        \draw(1.6, -1) -- (1.8, -1);
        \draw(1.6, -1.1) -- (1.8, -1.1);
        \draw(1.6, -1.2) -- (1.8, -1.2);
    \end{tikzpicture}    
\end{wrapfigure}
Si deve poter trovare un dispositivo che permetta di accumulare su di esso tanta carica 
ma che non abbia una differenza di potenziale troppo grande. Per poter aumentare la capacità 
di un oggetto che accumula carica è diminuire l'area dell'integrale del campo elettrico
del dispositivo. Un modo per farlo è potare lo zero del campo molto vicino alla superficie della sfera.
Per farlo si potrebbe mettere un'altra sfera intorno alla sfera iniziale 
di raggio $R_2 > R$ anche se si dovrebbe portarla a potenziale nullo (per esempio caricandola con $-Q$). 
Si sfrutta dunque il principio dell'\textbf{induzione elettrostatica}: dal punto di vista tecnico non ci si preoccupa 
di trovare un altro generaetore di carica e di potrare la sfera esterna a carica $-Q$, ma si effettua la 
\textbf{messa a terra}: dentro alla sfera esterna si ha campo non nullo e degli elettroni
fluiscano dalla Terra per annullare il campo generato dalla sfera interna. Il fenomeno dell'induzione elettrostatica 
dice che la sfera esterna e la Terra hanno lo stesso potenziale (neutre). In questo modo le cariche esterna eguagliano 
in modulo quelle interne: il campo elettrico si ha dunque solo all'esterno per il teorema di 
Gauss e si è realizzato un \textbf{condensatore}. Le due sfere prendono il nome di \textbf{armature} 
del condensatore e presenta. Allora il pontenziale diventa 
\begin{gather*}
    V(R_2) - V(R_1) = - \int_{R_1}^{R_2} E dr = - \frac{Q}{4\pi\epsilon_0} \left(\frac{1}{R_1} - \frac{1}{R_2}\right) 
\end{gather*}
Il valore specifico di una o l'altra armatura non è rilevante ma si può solo determinare la differenza di pontenziale 
tra di esse. Nel condensatore si può dunque scrivere
\begin{gather*}
    \Delta V = \frac{Q}{\mathcal{C}} \ \Longrightarrow \ \mathcal{C} = \frac{4\pi\epsilon_0}{R_2 - R_1}(R_1R_2) 
\end{gather*}
Nel limite in cui $R_1, R_2 >> R_2 - R_1$, (ossia nel limite di una piccola separazione tra le sfere),
si può chiamare $R_2 - R_1 = \delta$.  Dunque 
\begin{gather*}
    \mathcal{C} = 4\pi\epsilon_0R\frac{R}{\delta}
\end{gather*}

\subsection{Il condensatore a facce piane}
\begin{wrapfigure}{r}{0.4\textwidth}
    \centering
    \caption{}
    \begin{tikzpicture}
        \draw(0, -2) -- (0, 2) node[at end, left] {$Q$};
        \draw(2, -2) -- (2, 2) node[at end, above] {$V = 0$};
    \end{tikzpicture}    
\end{wrapfigure}
Utilizzare un condensatore sferico è inutilizzabile a livello pratico (anche
se è perfetto poiché permette di ottenere l'\textbf{induzione completa}) 
poiché si devono collegare entrambe le sfere del condensatore per poter usare la sua 
carica. Supponendo di avere due piani infiniti sui quali uno c'è carica $Q$, mentre l'altro 
è messo a masssa. Il teorema di unicità mi permette di definire il potenziale 
della faccia messa a terra (è nullo). Applicando Poisson si dimostra che il potenziale 
sulla faccia messa a terra è nulla (nel limite in cui i piani siano infiniti). Se i piani sono 
infiniti però si ha \textbf{induzione completa} e dunque il potenziale fuori dai 
piani è nullo. La carica che risiede sul piano messo a terra è esattamente $-Q$. \\
Il campo elettrico del primo piano che si trova a $x = 0$ è dato da
\begin{gather*}
    \begin{tikzpicture}
        \draw(0, 0) -- (5, 0);
        \draw(0, -1) -- (2.5, -1) -- (2.5, 1) -- (5, 1);
        \node at (-0.5, -1) {$-\dfrac{\sigma}{2\epsilon_0}$};
        \node at (5.3, 1) {$\dfrac{\sigma}{2\epsilon_0}$};
    \end{tikzpicture}
\end{gather*}
Il secondo piano ha un campo che invece è dato da
\begin{gather*}
    \begin{tikzpicture}
        \draw(0, 0) -- (5, 0);
        \draw(0, 1) -- (2.5, 1) -- (2.5, -1) -- (5, -1);
        \node at (5.5, -1) {$-\dfrac{\sigma}{2\epsilon_0}$};
        \node at (-0.5, 1) {$\dfrac{\sigma}{2\epsilon_0}$};
    \end{tikzpicture}
\end{gather*} 
Sommando i due campi si ha che
\begin{align*}
        x > d \ \Longrightarrow& \ E = 0 \\
        x < 0 \ \Longrightarrow& \ E = 0 \\
        0 < x < d \ \Longrightarrow& \ E = \frac{\sigma}{\epsilon_0}
\end{align*}
La differenza di potenziale invece è data dalla seguente
\begin{gather*}
    \Delta V = \frac{\sigma}{\epsilon_0} d = \frac{Q}{\epsilon_0 S} d = \frac{Q}{\mathcal{C}} \ \Longrightarrow \ \mathcal{C} = \frac{\epsilon_0 S}{d}
\end{gather*}

\section{Schermo elettrostatico}
Supponendo di avere un conduttore con forma qualunque e supponendo di avere delle
cariche all'esterno. Quello che si sa dei conduttori è che la carica si distribuisce
sulla superficie e, all'equilibrio si deve avere che il campo interno è nullo. La
$\sigma$ che si distribuisce sul conduttore non sarà la stessa ed il suo integrale
è nullo poiché il conduttore è neutro e dunque prende il nome di \textbf{condizione di neutralità 
del conduttore}. Supponendo ora che il conduttore abbia una cavita all'interno di esso,
le stessa cariche esterne, sulla superficie esterne si distribuiscono le cariche 
con una certa distribuzione $\sigma$ in modo tale da avere carica nulla. All'interno 
della cavità invece, non avendo cariche $\rho = 0$, $\nabla^{2} V = 0$ e inoltre, il bordo 
della cavità, essendo metallico ed essendo in statica, avrò un potenziale ben definito e costante. 
Utilizzando il teorema di unicità, si può dire che  il potenziale (soddisfacendo 
la condizione di Poisson) ha espressione:
\begin{gather*}
    V(r) = V_S
\end{gather*}
Questo vuol dire due cose: la cavità è \textbf{schermata} dai campi esterni a prescindere
dalla carica esterna al conduttore. La seconda cosa è che sulla superficie interna 
della cavità si ha $\sigma' = 0$ poiché il campo è continuo e uguale a zero e, per il teorema 
di Coulomb, sulla superficie del metallo esterno è nullo, allora in ogni punto del metallo 
non c'è carica e dunque non c'è accumulo di carica sulla superficie esterna. In condizioni statiche quello che è dentro
al metallo non è influenzato dalla carica circostante.

\subsection{Conduttore sferico cavo}
Considerando una cavità all'interno di un conduttore sferico e, ponendo una carica dentro la cavità, 
si sa che sulla superficie interna esiste una $\sigma$ che annulla il campo elettrico fuori 
dalla cavità: dunque il campo dentro al conduttore è $\vv{E} = 0$. Sulla superficie interna ,
l'integrale su $\sigma$ è 
\begin{gather*}
    \int_{S_{I}}^{} \sigma \ dS = -Q
\end{gather*}
poiché si prende una superficie di Gauss che continene tutta la cavità 
ma non esce dal metallo, per il teorema di Gauss, dato che il campo è nullo dentro al conduttore,
allora il flusso è nullo. Se il flusso è nullo, allora la carica che contiene la sueprficie di Gauss è nulla, 
dunque l'integrale della carica superficiale della cavità deve essere necessariamente $-Q$ 
per far sì che si bilanci la carica dentro la cavità $Q$ per soddisfare Gauss.
Sulla superficie esterna del conduttore si devono valutare due condizioni:
\begin{enumerate}
    \item \textbf{Isolato}: se il conduttore è isolato (e dunque anche neutro), 
    allora sulla superficie esterno si deve avere $\sigma'$ il cui integrale è opposto 
    a quello di $\sigma$ in modo che la carica sia nulla:
    \begin{gather*}
        \int_{S_{I}}^{} \sigma \ dS + \int_{S_{E}}^{}  \sigma' \ dS = 0 \ \Longrightarrow \ \int_{S_{E}}^{} \sigma' \ dS = -Q
    \end{gather*}
    Sulla superficie esterna c'è allora una carica $Q$ che è opposta a quello che è presente 
    sulla superficie interna. Se si facesse Gauss allora il campo è radiale e sulla sfera
    si avrebbe 
    \begin{gather*}
        V = \frac{Q}{4\pi\epsilon_0} \frac{1}{r} \ \Longrightarrow \ \sigma' = \frac{Q}{4\pi R^{2}}
    \end{gather*}
    L'unica informazione che si ha del conduttore è che il campo esterno è generato dalla 
    superficie esterna e dipende dalla carica interna $Q$. 
    \item \textbf{Sfera a massa}: si porta dunque la sfera al potenziale nullo. Per le equazioni di Poisson 
    non si ha potenziale esterno, dunque non si ha campo e allora non si ha $\sigma$ sulla superficie esterna: allora 
    il conduttore si è caricato e non è più neturo. 
\end{enumerate}

\subsection{Avvicinare le facce di un condensatore}
Nel conduttore a a facce piane e parallele, se ci si mette molto vicino ad una 
delle armature, allora vale l'approssimazione del piano infinto fino a che non si arriva veramente ad una distanza
dal bordo di una distanza comparabile a $\delta$ (ossia la distanza tra i due piani). Lo stesso vale se si prendono due 
piani e li si mettano molto vicini tra loro in modo tale che l'approssimazione di piano infinito
vale per entrambi i piani fino a che non ci si avvicina troppo al bordo. Per calcolare la Capacità 
si può considerare come se un condensatore avesse facce finite: questa approssimazione può 
essere utilizzato per determinare la capacità di un condensatore finito.

\subsection{Condensatori in serie}
Se si ponesse due condensatori $C_1$ e $C_2$ in serie tra loro, l'armatura
due del primo condensatore è allo stesso potenziale dell'armatura uno del secondo condensatore.
La domanda è capire che capacità ha questo sistema di condensatori 
se sono all'interno di un circuito con differenza di potenziale $\Delta V$.
\begin{gather*}
    Q_1 = \mathcal{C}_1 \Delta V_1 \qquad Q_2 = \mathcal{C}_2 \Delta V_2
\end{gather*}
Per cui $\Delta V = \Delta V_1 + \Delta V_2$. Se si immaginasse che le due facce allo stesso 
potenziale fossero una lastra isolata (e quindi neutra) con un determinato potenziale, allora si otterrebbe 
che $Q_1 = Q_2$, allora
\begin{gather*}
    Q = \mathcal{C}_T \Delta V = \mathcal{C}_T (\Delta V_1 + \Delta V_2)
\end{gather*}
Si sa anche che 
\begin{gather*}
    Q = \mathcal{C}_1 \Delta V_1 = \mathcal{C}_2 \Delta V_2
\end{gather*}
Dunque la capacità totale sarà
\begin{gather*}
    \mathcal{C}_T = \left(\frac{1}{\mathcal{C}_1} + \frac{1}{\mathcal{C}_2}\right)^{-1}
\end{gather*}

\subsection{Condensatore in parallelo}
Nei condensatori in parallelo, il vincolo è che $\Delta V_1 = \Delta V_2 = \Delta V$, allora 
le facce opposte dei condensatori hanno stesso potenziale, anche se non è più vero che
$Q_1 = Q_2$, ma la carica totale è la somma delle cariche
\begin{gather*}
    Q_T = Q_1 + Q_2 = \mathcal{C}_1\Delta V + \mathcal{C}_2 \Delta V \ \Longrightarrow \ \mathcal{C}_T = C_1 + C_2
\end{gather*}

\subsection{Problema del buco}
    Se si prendesse una carica e la si ponesse su di una sfera conduttrice di raggio $R$, 
    su ogni elemento $dS$ corrisponde una carica $dQ = \sigma \ dS$. Su ogni elemento di carica
    esiste una forza che spinge le cariche in avanti che si vuole calcolare. Questa forza 
    \begin{gather*}
        \vv{F_{dS}} = dQ \vv{E} 
    \end{gather*}
    Per il teorema di Coulomb si ha che
    \begin{gather*}
        dQ \vv{E} = \sigma dS \frac{\sigma}{\epsilon_0} \hat{u}  = \frac{\sigma^{2}}{\epsilon_0} dS \hat{u} 
    \end{gather*}
    l si ottiene una forza di pressione
    \begin{gather*}
        P = \frac{\left| F_u \right| }{dS} = \frac{\sigma^{2}}{\sigma}
    \end{gather*}
    Il problema è che al denominatore ci dovrebbe essere un due: questo perché la forza che risente
    ogni carica sul conduttore non può fare forza su di essa. Si deve valutare il campo su $dS$ di 
    tutto il resto della sfera:
    \begin{gather*}
        S = ds + \Sigma
    \end{gather*}
    Dove $\Sigma$ è la sfera senza il piccolo pezzo $dS$ (poiché $dS$ non genera forza su sé stesso). Il campo 
    generato dalla superficie
    \begin{gather*}
        \vv{E_S} = \frac{\sigma}{\epsilon_0}\hat{r} = \vv{E_{dS}} + \vv{E_\Sigma}   
    \end{gather*}
    Dove
    \begin{gather*}
        E_{dS} = \left\{\begin{array}{l}
            -\dfrac{\sigma}{2\epsilon_0} \hat{r} \quad r = R^{-} \\
            \dfrac{\sigma}{2\epsilon_0}\hat{r} \qquad r = R^{+}  
        \end{array}\right. \qquad E_S = \left\{\begin{array}{l}
            0 \qquad\quad r = R^{-} \\
            \dfrac{\sigma}{\epsilon_0 }\hat{r} \quad r = R^{+} 
        \end{array}\right.
    \end{gather*}
    Allora
    \begin{gather*}
        \vv{E}_\Sigma = \left\{\begin{array}{l}
            \dfrac{\sigma}{2\epsilon_0} \qquad r = R^{-} \\
            \dfrac{\sigma}{2\epsilon_0} \qquad r = R^{+}
        \end{array}\right. 
    \end{gather*}
    Il campo della sfera $\Sigma$ non è continuo, ma in un punto molto vicino al 
    buco è continuo. Quindi per la forza si deve scrivere che 
    \begin{gather*}
        \vv{F} = dQ\vv{E_\Sigma}  
    \end{gather*}
    Gli elettroni non possono dunque attraversare il buco poiché costituisce un vincolo
    geometrico (in realtà è quanto-meccanico). Quindi  pressione finale è 
    \begin{align}
        P = \frac{\sigma^{2}}{2\epsilon_0}
    \end{align}
    Poiché $\vv{E_\Sigma}$ presenta il due al denominatore. 

\subsection{Problema della carica immagine}
Preso un piano infinito messo a massa ed una carica $q$ posta a distanza $d$ dal piano.
Si deve dunque calcolare il campo elettrico in tutte le regioni dello spazio e la carica superficiale.
Se si dividesse lo spazio in regione destra e sinistra (II e I ), si vuole determinare 
il campo elettrico in entrambe le regioni. Il campo per I è nullo poiché si potrebbe vedere il piano infinto come 
    se fosse una sfera infinita (dunque se la sfera è a terra la carica interna è nulla).
    In questo caso, se il piano è messo a massa, per schermare il piano 
    si deve avere una carica negativa che viene dalla terra per bilanciare la carica nella regione II. 
    Per poter determinare il campo si risolve un secondo problema nel quale il piano a terra presenta due 
cariche $-q$ e $q$ a distanza $d$. Per verificare che questi due problemi diano la stessa 
soluzione si sfrutta il teorema di unicità: l'equazione di Poisson mi permette di dire che 
$\rho_1 = \rho_2$ ossia la $\rho$ dei due problemi. Sul piano che taglia le due cariche,
il potenziale che generano le cariche è
\begin{gather*}
    V_q = \frac{q}{4\pi\epsilon_0} \frac{1}{r}  \qquad V_{-q} = -\frac{q}{4\pi\epsilon_0} \frac{1}{r}
\end{gather*}
In tutti i punti del piano i potenziali si equilibrano e dunque il potenziale è nullo.
Allora nel primo problema $V_{II}(\infty ) = 0$. La densità di carica è indotta in modo tale che 
nella parte I il campo sia nulla. 
QUesto problema prende il nome di problema della carica 
immagine perché il metallo è come se fosse uno specchio (posso mettere nella regione 
che si è risolta (la I) una carica (o più cariche) specchiata rispetto a $q$). Il campo sarà dunque 
la somma dei campi delle cariche. Per determinare invece $\sigma$ sul piano, si sa che 
\begin{gather*}
    E(x = 0^{+}) = \frac{\sigma}{\epsilon_0} \hat{x} 
\end{gather*}
Nel "problema immagine",si pu determinare il campo sulla superficie, dunque il campo 
$E(x = 0^{+})$ si può scrivere come il campo di due cariche uniformi 
\begin{gather*}
    E(x = 0^{+}) = \frac{q}{4\pi\epsilon_0} \frac{d\hat{x} + r \hat{r}  }{(d + r^{2})^{3/2}} - \frac{q}{4\pi\epsilon_0} \frac{-d\hat{x} - r\hat r  }{(d^{2} + r^{4})^{3/2}} = -\frac{2qd}{4\pi\epsilon_0 (d^{2} + r^{2})^{3/2}}
\end{gather*}


\end{document}